% GENERATED from content/data/enrichment.json by mxbuild.enrichment; do not edit.
\mxunit{enrichment-bayes-notes}{Inverser une probabilité conditionnelle}
\mxsubtitle{Première--Terminale · probabilités \textperiodcentered{} Leçon et méthode \textperiodcentered{} 07}

\begin{prerequisites}
Arbres pondérés, intersections et probabilités conditionnelles.
\end{prerequisites}

\begin{keypoint}[Idées et méthode]
$P(A|B)$ désigne la probabilité de A lorsqu’on sait B réalisé. Si $P(B) > 0$, $P(A|B) = \frac{P(A \cap  B)}{P(B)}$. L’information B restreint l’univers considéré.

$P(A \cap  B) = P(A)P(B|A)$. Si A et son contraire forment les deux premières branches d’un arbre, $P(B) = P(A)P(B|A) + P(\text{non }A)P(B|\text{non }A)$.

La formule de Bayes résulte de ces deux égalités. $P(A|B)$ n’est généralement pas égal à $P(B|A)$ : les dénominateurs correspondent à des populations différentes.
\end{keypoint}

\begin{example}
On choisit A avec probabilité $\frac{1}{3}$ et B avec probabilité $\frac{2}{3}$. L’urne A contient 3 rouges et 1 bleue ; B contient 1 rouge et 3 bleues. $P(\text{rouge}) = (\frac{1}{3})(\frac{3}{4}) + (\frac{2}{3})(\frac{1}{4}) = \frac{5}{12}$. Sachant que la boule est rouge, $P(A|\text{rouge}) = \frac{\frac{1}{4}}{\frac{5}{12}} = \frac{3}{5}$.
\end{example}

\begin{figurebox}[Seules les branches conduisant à une boule rouge sont représentées.]
\begin{tikzpicture}[x=1.6cm,y=.8cm]\draw[mxmint,thick](0,0)--(1,1)node[midway,above]{$\frac13$};\draw[mxmint,thick](0,0)--(1,-1)node[midway,below]{$\frac23$};\node at(1.2,1){$A$};\node at(1.2,-1){$B$};\draw[mxink,thick](1.5,1)--(3,1.5)node[midway,above]{$\frac34$};\draw[mxink,thick](1.5,-1)--(3,-.5)node[midway,above]{$\frac14$};\node[right]at(3,1.5){$A\cap R:\ \frac14$};\node[right]at(3,-.5){$B\cap R:\ \frac16$};\end{tikzpicture}
\end{figurebox}

\begin{activity}[À explorer]
Représenter l’exemple sur 120 tirages fictifs respectant exactement les proportions. Compter les boules rouges venant de chaque urne.
\end{activity}

